% !TeX root = main.tex

% ============================================================
% thesis-sync entry: 2026-07-19_inchworm-gui
% Type: design decision (software architecture) — reinstates the paper's
%       commented-out "Pneumatic System Graphical User Interface" subsection.
% Part of the INCHWORM ROBOT block (prior work, precedes the lizard sections).
% NOTE: this is the GUI of the eight-channel *solenoid* pneumatic system
% (Section~\ref{sec:inchworm-pneumatics}); it is distinct from the manual
% syringe-pump control interface documented for the lizard robot.
% Insertable section block (no preamble). Requires: graphicx.
% ============================================================

\section{GRAPHICAL USER INTERFACE}
\label{sec:inchworm-gui}

To orchestrate complex gait sequences and to facilitate intuitive operation
of the soft robot, a custom Graphical User Interface (GUI) was developed to
establish real-time, bidirectional communication with the pneumatic control
hardware. The interface is structurally divided into two primary
environments, to separate trajectory generation from execution.

The \textit{Program} tab (Fig.~\ref{fig:inchworm-programtab}) provides a
dedicated workspace in which the user designs and preprograms specific
robotic ``actions.'' An action is characterized by a synchronized set of
target-pressure trajectories across all active pneumatic channels. These
continuous trajectories are constructed by interpolating between discrete
$(t, P)$ keypoints---representing time and target pressure,
respectively---defined over a finite duration. This allows for precise
temporal control over the inflation and deflation dynamics of the individual
actuator segments.

\begin{figure}[tbh]
 \centering
  \includegraphics[height=38mm]{figures/programTab.eps}
  \vspace*{-2mm}
  \caption{Program tab of the graphical user interface, in which the user
  preprograms the robot's actions by defining target-pressure curves through
  $(t,P)$ keypoints.}
  \label{fig:inchworm-programtab}
\end{figure}

Once formulated, these multi-channel actions are serialized and stored as
JSON files within a centralized action library. This library is directly
accessible from the \textit{Live Control} tab
(Fig.~\ref{fig:inchworm-controltab}). Furthermore, the \textit{Live Control}
tab features real-time plotting capabilities, continuously overlaying the
target-pressure curves with the measured sensor outputs for each channel.
This provides immediate visual feedback on the controller's tracking
performance, and the resulting graphical data can be exported, capturing up
to 60 seconds of plot history for subsequent empirical analysis. In this
deployment environment, the user can plan and execute combinations of these
preprogrammed actions, evaluating the soft robot's macroscopic locomotion and
behavioral responses.

\begin{figure}[tbh]
 \centering
  \includegraphics[height=38mm]{figures/controltab.eps}
  \vspace*{-2mm}
  \caption{Live Control tab of the graphical user interface, from which the
  user launches sequences of preprogrammed actions to control the robot in
  real time while monitoring target-versus-measured pressure tracking.}
  \label{fig:inchworm-controltab}
\end{figure}
